
\documentclass[11pt]{article}
\usepackage[margin=1in]{geometry}

% Core packages
\usepackage{amsmath,amssymb}
\usepackage{tikz-cd}
\usepackage{multicol}

% Paragraphs
\setlength{\parindent}{0pt}
\setlength{\parskip}{1\baselineskip}

\title{Log de Trabajo}
\author{Jaime Barrera Ortiz}
\date{\today}

\begin{document}
\maketitle

\section{Decisión del Preprocesamiento}


\subsection{Diagnóstico Inicial}

El análisis exploratorio mostró que los CETEs presentan observaciones mensuales constantes, mientras que los MBONOs (especialmente los de 20Y y 30Y) contienen ausencias significativas (hasta 30$\%$ de meses sin subasta). No existen errores de captura, pero la matriz de datos es incompleta.

\subsection{Métodos de Imputación}\label{subsec:metodos-imputacion}

Se descartó la eliminación estricta debido a que resulta en 0 meses con observaciones completas para los 9 tenores simultáneamente. Forward Fill se adoptó como la mejor decisión debido a que preserva la continuidad temporal de la serie, ademas sus limitaciones son bien conocidas y se mencionan a continuación.

\subsection{Riesgos Identificados}

Al trabajar con ffill se introduce una correlación espuria, afectando principalmente a la PC1 (Nivel) ya que puede inflar artificialmente la varianza si los periodos coinciden con baja volatilidad real, y a la PC3 (Curvatura) volviéndola vulnerable por su baja sincronización entre 20Y y 30Y, una alta correlación puede reflejar el calendario administrativo de banxico más que un factor de riesgo real.

\subsection{Justificación}

Dado que el objetivo es identificar estructura en la curva completa, ffill permite mantener los 9 instrumentos en el modelo, manteniendo cierta discreción al interpretar los componentes principales, especialmente la PC3.


\section{Modelado de los Datos}

\subsection{Análisis de Componentes Principales (PCA)}

Se seleccionó el PCA como herramienta central para responder la primera pregunta del ejercicio (\textit{¿Existe estructura en cómo se mueven juntas las 9 maturities?}) por las siguientes razones :

\begin{enumerate}
    \item \textbf{Reducción de Dimensionalidad:} La curva de tasas mexicana posee 9 dimensiones observadas. El PCA permite descomponer la varianza total en factores ortogonales (Nivel, Pendiente, Curvatura), transformando un problema de varias dimensiones en uno de baja dimensión. 
    
    \item \textbf{Compatibilidad con la Serie:} Dado que se adoptó \texttt{ffill} para preservar la continuidad temporal, como se explica en la Subsección~\ref{subsec:metodos-imputacion}, el PCA sobre la matriz de correlaciones estandarizadas es matemáticamente consistente. Al trabajar sobre cambios mensuales, se mitiga la no-estacionariedad inherente a los niveles de tasa, asegurando que los componentes reflejen dinámicas de corto plazo y no tendencias deterministas.
    
    \item \textbf{Interpretabilidad:} Aunque \texttt{ffill} puede inflar la correlación espuria (especialmente en PC3 debido a la asincronía 20Y/30Y), el PCA sigue siendo el método estándar en literatura de curvas de rendimiento (Litterman-Scheinkman, 1991). Si los tres primeros componentes explican $>95\%$ de la varianza, confirmamos que la curva mexicana se comporta de manera similar a mercados desarrollados, validando la existencia de factores robustos incluso bajo imputación conservadora.
    
    \item \textbf{Fundamento para Forecasting (Preguntas 2 y 3):} Los scores del PCA proporcionan series temporales estacionarias y de baja dimensionalidad que pueden utilizarse como features en modelos de forecasting. Generando un puente entre la identificación de estructura y la mejora de pronósticos, 

\end{enumerate}

\noindent \textbf{Nota:} Aunque existen métodos alternativos como PCA probabilístico que podrían manejar mejor los datos faltantes, sin embargo, dada la inviabilidad de la eliminación estricta y la necesidad de mantener comparabilidad con la literatura clásica de curvas de tasas, el PCA sobre datos ffill representa el equilibrio  entre rigor estadístico y viabilidad computacional.


\pagebreak

\section{Forecasting de la Curva de Rendimientos}

\subsection{Estrategia de Modelado}

Para responder a las preguntas 2 y 3 del ejercicio (pronóstico de la curva a horizontes de 1 trimestre y 4 trimestres), se adoptó una estrategia en dos etapas: primero se pronostican las trayectorias temporales de los tres factores latentes principales obtenidos del PCA (Nivel, Pendiente y Curvatura), y posteriormente se reconstruye la estructura temporal completa proyectando dichos factores a través de la matriz de cargas factoriales ($V_{1:3}$), desestandarizando a unidades de tasa de interés real ($\%$).

Se implementaron tres modelos con fundamentos conceptuales distintos:

\begin{enumerate}
    \item \textbf{Random Walk (Benchmark ingenuo):} Asume que la curva futura esperada condicional a la información presente es idéntica a la última curva observada ($\mathbb{E}[S_{t+h} \mid \mathcal{F}_t] = S_t$). En econometría de tasas de interés, este modelo funge como la hipótesis nula obligatoria: cualquier formulación dinámica debe probar empíricamente que reduce el error cuadrático respecto a la curva estática.
    
    \item \textbf{Ornstein-Uhlenbeck / AR(1) desacoplado (Baseline de reversión a la media):} Modela cada componente principal de forma independiente como un proceso de difusión estocástica con atracción hacia un nivel medio de largo plazo. Físicamente representa un oscilador sobreamortiguado sujeto a perturbaciones brownianas; financieramente refleja la postura de política monetaria de Banxico, donde desviaciones severas en el nivel o la pendiente tienden a revertir gradualmente hacia un ancla macroeconómica.
    
    \item \textbf{Vector Autoregression - VAR(p):} Permite la retroalimentación dinámica y el acoplamiento cruzado entre los tres modos normales. A diferencia del modelo desacoplado, un VAR captura cómo un shock en la pendiente (ej. inversión de curva por alza de tasas de referencia) impacta rezagadamente al nivel general o a la curvatura en meses subsecuentes. El orden de rezagos $p$ se determina de forma adaptativa mediante el Criterio de Información de Akaike (AIC) para balancear bondad de ajuste y parsimonia.
\end{enumerate}

\subsection{Opciones Descartadas y Justificación}

Durante el diseño de la etapa predictiva se evaluaron y descartaron las siguientes alternativas:

\begin{itemize}
    \item \textbf{Ajuste directo de series sobre los 9 instrumentos (VAR de dimensión 9):} Se descartó debido a que estimar un sistema multivariado sobre las 9 maturities simultáneamente requiere estimar $9^2 \cdot p$ parámetros. Con un historial mensual de alrededor de 240 observaciones, esto conduce de inmediato a la maldición de la dimensionalidad, sobreajuste severo e inestabilidad numérica producto de la multicolinealidad extrema (correlación $>0.95$ entre tenores contiguos). El modelado en el subespacio latente ortogonal de dimensión 3 resuelve este problema por construcción.
    
    \item \textbf{Modelos de Aprendizaje Profundo / Machine Learning (LSTM, MLP, XGBoost):} Se descartaron porque requieren miles de observaciones independientes para aprender representaciones sin memorizar el ruido estocástico. En series macroeconómicas mensuales de muestra moderada, estas arquitecturas tienden a sobreajustar agresivamente y carecen de una formulación analítica cerrada para la propagación de covarianzas de error.
    
    \item \textbf{Modelos Paramétricos de Nelson-Siegel Dinámico:} Aunque son habituales en la práctica de bancos centrales, se descartó su calibración debido a que optimizan funciones no lineales con parámetros de decaimiento ($\lambda$) que presentan mínimos locales y volatilidad numérica en periodos donde los MBONOs largos tuvieron subastas ausentes. Los autovectores del PCA ofrecen una base modal fija, lineal y ortogonal que evita algoritmos de optimización no lineal inestables.
\end{itemize}

\subsection{Prevención de Data Leakage en la Evaluación Fuera de Muestra}

Un riesgo crítico identificado en el diseño original fue el \textit{look-ahead bias} en la validación temporal. Si se ajusta el PCA sobre la muestra completa y luego se evalúa el pronóstico sobre los últimos meses de scores, los autovectores y los parámetros de estandarización ($\mu, \sigma$) ya habrían incorporado la estructura de covarianza del periodo de prueba.

Para garantizar una evaluación estrictamente fuera de muestra (\textit{Out-of-Sample}):
\begin{enumerate}
    \item Se implementó un esquema de evaluación por ventana móvil expansiva (\textit{walk-forward rolling origin}).
    \item En cada punto de corte temporal $t_{\text{cutoff}}$, se reentrena el PCA y se recalculan la media y desviación estándar empleando exclusivamente los datos observados hasta ese mes ($t \le t_{\text{cutoff}}$).
    \item El pronóstico proyectado se desestandariza hacia el espacio de tasas reales y se compara a ciegas contra las tasas de subasta efectivamente observadas en el mercado en los $h$ meses posteriores.
\end{enumerate}

\subsection{Propagación de Incertidumbre y Densidad Predictiva (Fan Chart)}

Dado que los pronósticos puntuales no transmiten la dispersión natural del riesgo a través del horizonte temporal, se incorporó la estimación explícita de densidad de probabilidad:

\begin{enumerate}
    \item \textbf{Evolución temporal de la varianza:} A partir de la matriz de covarianza de los residuos del VAR o de los errores de innovación en AR(1) y Random Walk, se calcula la matriz de covarianza del error de pronóstico $\Omega(h)$ para cada horizonte $h \in \{1, 3, 6, 12\}$.
    \item \textbf{Proyección a cada tenor:} Mediante la forma cuadrática $\Sigma_{\text{curve}}(h) = V_{1:3}^T \Omega(h) V_{1:3}$, se obtiene la varianza analítica de cada instrumento en unidades estandarizadas y, tras multiplicar por la desviación estándar del instrumento, en puntos porcentuales de tasa real.
    \item \textbf{Visualización integrada (Fan Chart):} Se descartó presentar densidades estáticas aisladas en favor de un gráfico de abanico (\textit{Fan Chart}) continuo. Este conecta la serie histórica de los últimos meses observados del instrumento con la trayectoria proyectada, envuelta en bandas de probabilidad del $50\%$, $80\%$ y $95\%$, ilustrando cómo el cono de incertidumbre se ensancha conforme aumenta el plazo de predicción.
\end{enumerate}

\end{document}
